% ═══════════════════════════════════════════════════════════════════════════
% VERSUCHSPROTOKOLL: Die Federspannkraft (UE 4)
% Physik Klasse 7 – Kräfte – 2 Seiten
% ═══════════════════════════════════════════════════════════════════════════

\documentclass[10pt,a4paper]{article}

\usepackage[utf8]{inputenc}
\usepackage[T1]{fontenc}
\usepackage[ngerman]{babel}
\usepackage{helvet}
\renewcommand{\familydefault}{\sfdefault}

\usepackage[margin=15mm,top=12mm,bottom=12mm]{geometry}
\usepackage{xcolor}
\usepackage{tikz}
\usepackage{array}
\usepackage{amsmath}

% ═══════════════════════════════════════════════════════════════════════════
% FARBEN
% ═══════════════════════════════════════════════════════════════════════════
\definecolor{physik}{HTML}{2c5aa0}
\definecolor{grau}{HTML}{888888}
\definecolor{hellgrau}{HTML}{e8e8e8}
\definecolor{raster}{HTML}{c0c0c0}

% ═══════════════════════════════════════════════════════════════════════════
% BEFEHLE
% ═══════════════════════════════════════════════════════════════════════════
\newcommand{\luecke}[1]{\underline{\hspace{#1}}}
\newcommand{\punkte}[1]{\hfill\textcolor{grau}{\small/#1}}
\newcommand{\stern}[1]{\textcolor{physik}{\textbf{#1}}}
\newcommand{\fv}[1]{\textit{#1}}
\newcommand{\op}[1]{\textbf{#1}}

% Platz für Rechnung (3 Zeilen à 8mm)
\newcommand{\rechenplatz}{\vspace{24mm}}

% ═══════════════════════════════════════════════════════════════════════════
% LAYOUT
% ═══════════════════════════════════════════════════════════════════════════
\pagestyle{empty}
\setlength{\parindent}{0pt}
\setlength{\parskip}{0pt}

% ═══════════════════════════════════════════════════════════════════════════
\begin{document}

% ───────────────────────────────────────────────────────────────────────────
% KOPFBEREICH
% ───────────────────────────────────────────────────────────────────────────
\begin{tikzpicture}[remember picture, overlay]
    \fill[physik] (current page.north west) rectangle ([yshift=-10mm]current page.north east);
    \node[anchor=west, white, font=\normalsize\bfseries] at ([xshift=15mm, yshift=-5mm]current page.north west)
        {Versuchsprotokoll: Die Federspannkraft};
    \node[anchor=east, white, font=\small] at ([xshift=-15mm, yshift=-5mm]current page.north east)
        {Physik 7};
\end{tikzpicture}

\vspace{5mm}

{\small Name: \luecke{4cm} \hfill Klasse: \luecke{1cm} \hfill Datum: \luecke{2cm}}

\vspace{2mm}
\hrule height 0.3pt
\vspace{3mm}

% ───────────────────────────────────────────────────────────────────────────
% FRAGESTELLUNG + DURCHFÜHRUNG + MESSWERTE (kompakt)
% ───────────────────────────────────────────────────────────────────────────
\textbf{\textcolor{physik}{Fragestellung:}} Wie hängt die Dehnung einer Feder von der angreifenden Kraft ab?

\vspace{2mm}

\textbf{\textcolor{physik}{Durchführung:}} Öffne die Simulation. Wähle Feder 1. Hänge Gewichte an und lies die Dehnung ab.

\vspace{3mm}

\textbf{\textcolor{physik}{Messwerte:}} \punkte{5}

\vspace{2mm}

\begin{center}
\begin{tabular}{|c|c|c|c|c|c|}
\hline
\rowcolor{hellgrau}
\textbf{Kraft \fv{F} in N} & 1 & 2 & 3 & 4 & 5 \\
\hline
\textbf{Dehnung \fv{s} in cm} & & & & & \\[6mm]
\hline
\end{tabular}
\end{center}

\vspace{3mm}

% ───────────────────────────────────────────────────────────────────────────
% DIAGRAMM (groß, 5mm Kästchenpapier)
% ───────────────────────────────────────────────────────────────────────────
\textbf{\textcolor{physik}{Diagramm:}} \op{Zeichne} ein \fv{F}-\fv{s}-Diagramm mit Ausgleichsgerade. \punkte{6}

\vspace{2mm}

\begin{center}
\begin{tikzpicture}
    % 5mm Kästchenpapier
    \draw[hellgrau, very thin] (0,0) grid[step=0.5] (13,6.5);
    \draw[raster, thin] (0,0) grid[step=1] (13,6.5);

    % Achsen
    \draw[->, very thick] (0,0) -- (13.5,0) node[right] {\fv{s} in cm};
    \draw[->, very thick] (0,0) -- (0,7) node[above] {\fv{F} in N};

    % Beschriftung x-Achse
    \foreach \x in {0,1,2,3,4,5,6,7,8,9,10,11,12,13} {
        \draw[thick] (\x,-0.1) -- (\x,0.1);
        \node[below] at (\x,-0.15) {\small\x};
    }

    % Beschriftung y-Achse
    \foreach \y in {0,1,2,3,4,5,6} {
        \draw[thick] (-0.1,\y) -- (0.1,\y);
        \node[left] at (-0.15,\y) {\small\y};
    }
\end{tikzpicture}
\end{center}

\vspace{2mm}

% ───────────────────────────────────────────────────────────────────────────
% BEOBACHTUNG + ERKENNTNIS
% ───────────────────────────────────────────────────────────────────────────
\textbf{\textcolor{physik}{Beobachtung:}} \op{Beschreibe}, was du im Diagramm erkennst. \punkte{3}

\vspace{1mm}

Wenn die Kraft verdoppelt wird, dann \luecke{4.5cm}. Die Punkte liegen auf einer \luecke{2.5cm}.

\vspace{3mm}

\textbf{\textcolor{physik}{Erkenntnis (Hooke'sches Gesetz):}} \punkte{5}

\vspace{1mm}

\begin{minipage}[t]{0.48\textwidth}
Die Dehnung \fv{s} ist \luecke{3cm} zur Kraft \fv{F}.

\vspace{3mm}

\textbf{Formel:} \fbox{\;\Large $\fv{F} = $\luecke{1.5cm}\;\;}
\end{minipage}
\hfill
\begin{minipage}[t]{0.48\textwidth}
\op{Gib an}, was die Federkonstante \fv{D} bedeutet:

\vspace{8mm}

Einheit: $[\fv{D}] = $ \luecke{2cm}
\end{minipage}

% ═══════════════════════════════════════════════════════════════════════════
% SEITE 2
% ═══════════════════════════════════════════════════════════════════════════
\newpage

\textbf{\textcolor{physik}{Aufgaben \stern{★}}}

\vspace{3mm}

\textbf{1.} \op{Ermittle} aus deinem Diagramm: Welche Kraft dehnt die Feder um 6\,cm? \punkte{2}

\vspace{6mm}

\textbf{2.} Bei $\fv{F} = 2$\,N dehnt sich eine Feder um 4\,cm. \op{Bestimme} die Dehnung bei 4\,N und \op{begründe}. \punkte{2}

\vspace{10mm}

\textbf{3.} \op{Nenne} zwei Alltagsgegenstände mit Federn. \op{Erkläre} bei einem die Funktion der Feder. \punkte{2}

\vspace{12mm}

\hrule height 0.2pt
\vspace{4mm}

\textbf{\textcolor{physik}{Aufgaben \stern{★★}}}

\vspace{3mm}

\textbf{4.} Eine Feder hat $\fv{D} = 0{,}5\;\frac{\text{N}}{\text{cm}}$. \op{Berechne} die Kraft für eine Dehnung von 12\,cm. \punkte{3}

\rechenplatz

\textbf{5.} \op{Berechne} die Kraft, um eine Feder mit $\fv{D} = 2\;\frac{\text{N}}{\text{cm}}$ um 5\,cm zu dehnen. \punkte{3}

\rechenplatz

\hrule height 0.2pt
\vspace{4mm}

\textbf{\textcolor{physik}{Aufgaben \stern{★★★}}}

\vspace{3mm}

\textbf{6.} \op{Ermittle} die Federkonstante \fv{D} von Feder 1 aus deinem Diagramm. \op{Erkläre} ihre Bedeutung. \punkte{4}

\rechenplatz

\textbf{7.} Wechsle zu \textbf{Feder 2 (hart)}. \op{Bestimme} deren Federkonstante bei $\fv{F} = 3$\,N und \op{vergleiche}. \punkte{4}

\rechenplatz

\textbf{8.} \op{Berechne} \fv{D} von Feder 1 in $\frac{\text{N}}{\text{m}}$. \;\;{\small(Hinweis: $1\;\frac{\text{N}}{\text{cm}} = 100\;\frac{\text{N}}{\text{m}}$)} \punkte{2}

\vspace{10mm}

% ───────────────────────────────────────────────────────────────────────────
% BEWERTUNG
% ───────────────────────────────────────────────────────────────────────────
\vfill

\hrule height 0.3pt
\vspace{2mm}

{\small
\begin{tabular}{|l|l|c|c|}
\hline
\rowcolor{hellgrau}
\textbf{Niveau} & \textbf{Aufgaben} & \textbf{Max. BE} & \textbf{Erreicht} \\
\hline
\stern{★} Basis & Protokoll + Aufg. 1–3 & 25 & \\
\hline
\stern{★★} Standard & + Aufg. 4–5 & 31 & \\
\hline
\stern{★★★} Erweitert & + Aufg. 6–8 & 41 & \\
\hline
\end{tabular}
\hfill
\textcolor{grau}{\footnotesize Physik 7 – Kräfte – UE 4}
}

\end{document}
